\uuid{st4S}
\titre{Systèmes d'équations linéaires}
\auteur{hanani ; mzari}
\organisation{exo7}
\licence{CC BY-NC-SA 4.0}
\datecreate{2020-01-20}
\niveau{L1}
\module{Algèbre}
\chapitre{Déterminant, système linéaire}
\sousChapitre{Système linéaire, rang}
\difficulte{3}
\theme{Systèmes d'équations linéaires}
\exo7qcmid{lille-2.systemes.025}

\contenu{
\begin{question}
\qclassification{200.03}
\qtags{motcle=Système linéaire à $n$ inconnus et à paramètres}
On considère le système d'équations, d'inconnue $(x_1,x_2,\dots,x_n) \in \Rr^n$, $n \ge 2$ : et de paramètre des réels  $a,b$ :
$$(\mathtt{S})
\left\{\begin{array}{ccc}
x_1+x_2 + x_3+\dots \quad \dots+ x_n&=&1\\
ax_1+bx_2 + bx_3 +  \dots   +bx_n&=&1\\
ax_1+ax_2+bx_3 + \dots +bx_n&=&1\\ \vdots &\vdots&\vdots\\
ax_1+\dots + ax_{n-1}+bx_n&=&1,\end{array}\right.$$
où $a$ et $b$ sont des paramètre réels. Quelles sont les assertions vraies ?
\begin{answers}
\good{$(\mathtt{S}) \Leftrightarrow
\left\{\begin{array}{rcc}
x_1+x_2+x_3+\dots+ x_n&=&1\\
(b-a)[x_2+x_3+\dots+ x_n]&=&1-a\\
(b-a)[x_3+\dots +x_n]&=&1-a\\
\vdots &\vdots&\vdots\\
(b-a)x_n&=&1-a.\\
\end{array}\right.$}
\bad{Si $a=b$, $(\mathtt{S})$ admet une infinité de solutions.}
\bad{Si $a\neq b$, $(\mathtt{S})$ n'admet pas de solution.}
\good{$(\mathtt{S})$ admet une infinité de solutions si et seulement si $a=b=1$.}
\end{answers}
\begin{explanations} L'algorithme de Gauss donne :
$$(\mathtt{S}) \Leftrightarrow
\left\{\begin{array}{rcc}
x_1+x_2+x_3+\dots+ x_n&=&1\\
(b-a)[x_2+x_3+\dots+ x_n]&=&1-a\\
(b-a)[x_3+\dots +x_n]&=&1-a\\
\vdots &\vdots&\vdots\\
(b-a)x_n&=&1-a.
\end{array}\right.$$
\begin{enumerate}
\item[-] Si $a=b=1$, le système $(\mathtt{S})\Leftrightarrow x_1+x_2+\dots+ x_n=1$. Donc $(\mathtt{S})$ admet une infinité de solutions.
\item[-] Si $a=b \neq 1$, $(\mathtt{S})$ n'admet pas de solution.
\item[-] Si $a\neq b$, $(\mathtt{S})$ admet une unique  solution.
\end{enumerate}
\end{explanations}
\end{question}
}
